%% kapa: exam=1 round=35 year=2024 period=1 form=A subject=economics subject_ko=경제학원론 seq=16 number=56
%% source: Q-Net 감정평가사 기출문제 (cst00309) artlSeq=5213229 file=raw/1008/5213229/1교시 문제지 원본.pdf
%% conversion: kordoc 4.24.51 -> tools/gen1.py (auto); review=proofread-sonnet; flags=none
\begin{question}{56}
독점 기업 A의 비용 함수는 $C(Q)=750+5Q$이고, 역수요함수는 $P=140-Q$이다. 이 기업이 ‘독점을 규제하는 법률’에 따라 한계비용과 동일하게 가격을 설정한다면, 이에 관한 설명으로 옳은 것은? (단, $Q$는 수량, $P$는 가격이다.)
\begin{choices}
  \item 양(+)의 이윤을 얻는다.
  \item 이윤은 0이다.
  \item 손실이 375이다.
  \item 손실이 450이다.
  \item 손실이 750이다.
\end{choices}
\end{question}
\answer{5}
